% Large Language Models (LLMs) show promise for structured generation tasks like equation discovery, yet their outputs are highly sensitive to prompt phrasing---a phenomenon we term \textit{instruction brittleness}. We identify a core limitation: static prompts cannot adapt to the evolving state of a multi-step generation process, causing models to plateau at suboptimal solutions. To address this, we propose \ours, which reframes prompt engineering as a \textit{sequential decision problem}. Our method maintains a discrete library of reasoning strategies (e.g., ``simplify the expression,'' ``explore alternative structures'') and uses Bayesian Optimization to select the optimal instruction at each generation step based on numerical feedback. Experiments on PDE discovery benchmarks show that adaptive instruction selection significantly outperforms fixed prompts, achieving higher symbolic recovery rates with more parsimonious solutions.

\begin{abstract}
	Large Language Models (LLMs) show promise for equation discovery, yet their outputs are highly sensitive to prompt phrasing—a phenomenon we term \textit{instruction brittleness}. Static prompts cannot adapt to the evolving state of a multi-step generation process, causing models to plateau at suboptimal solutions. To address this, we propose \textsc{NeuroSym-BO}, which reframes prompt engineering as a \textit{sequential decision problem}. Our method maintains a discrete library of reasoning strategies and uses Bayesian Optimization to select the optimal instruction at each step based on numerical feedback. Experiments on PDE discovery benchmarks show that adaptive instruction selection significantly outperforms fixed prompts, achieving higher recovery rates with more parsimonious solutions.
\end{abstract}



